% Fragment autonome de preuve — ne contient pas de préambule.
% Casimir choisi positif sur les séries principales.
\section{Une borne spectrale pour le cercle complet}

Posons $G=\mathrm{SL}_2(\mathbb R)$, $\Gamma=\mathrm{SL}_2(\mathbb Z)$,
$K=\mathrm{SO}_2(\mathbb R)$ et $X=\Gamma\backslash G$, muni de la
probabilité de Haar $\mu$. Nous utilisons la représentation unitaire
$R(u)f(g)=f(gu)$ et les moyennes
\[
 \mathcal M_s f=\int_K f(\Gamma k a_s)\,dk,
 \qquad a_s=\operatorname{diag}(e^{-s},e^s),\quad s\ge0.
\]
Choisissons les trois champs invariants à gauche $X_0,Y_0,\Theta$ dans
la normalisation où
\[
 \mathcal D=-(X_0^2+Y_0^2+\Theta^2)=\mathcal C-2\Theta^2
\]
est positif, et où le Casimir $\mathcal C$ agit par le Laplacien
hyperbolique sur les fonctions invariantes à droite par $K$. Le générateur
compact satisfait $\Theta e_n=in e_n/2$ sur un vecteur du $K$-type $n$.
Nous prenons les réalisations autoadjointes naturelles et posons
\[
 W^r(X)=\operatorname{Dom}(1+\mathcal D)^{r/2},\qquad
 \|f\|_{W^r}=\|(1+\mathcal D)^{r/2}f\|_2.
\]
Le choix des signes de $\mathcal C$ est donc fixé par cette identité;
sa valeur propre dans une série principale est $\lambda\ge1/4$.

\begin{proposition}[Borne spectrale d'ordre $7/4$]
Il existe une constante $C$ telle que, pour $s\ge0$ et $f\in W^{7/4}(X)$,
\[
 |\mathcal M_s f-\mu(f)|
 \le C(1+s)e^{-s}\|f\|_{W^{7/4}}.
 \tag{S1}
\]
La moyenne du membre de gauche est prise avec la trace naturelle de $f$
sur la courbe $\Gamma K a_s$.
\end{proposition}

\begin{proof}
Écrivons $P_K=\int_K R(k)\,dk$. La moyenne recherchée est la valeur de
$P_KR(a_s)f$ au point $\Gamma e$, ou, comme cette fonction est
$K$-invariante, au point $i$ de la surface modulaire.

Les représentations non sphériques ne contribuent pas à $P_K$.
Dans le spectre sphérique de niveau~1, la représentation triviale donne
$\mu(f)$; toutes les autres représentations sont tempérées, en vertu de
l'absence de valeurs propres exceptionnelles pour le groupe modulaire.
On conserve à la fois le spectre cuspidal et le spectre continu.

Dans chaque représentation sphérique $\pi$, choisissons une base
orthonormale $(e_n)$ de $K$-types, avec $e_0$ sphérique, et écrivons
$f_\pi=\sum_n f_{\pi,n}e_n$. Chacun des deux vecteurs $e_n,e_0$ engendre
un espace de dimension~1 sous $K$. La borne de Cowling--Haagerup--Howe
s'applique donc \emph{séparément à chaque type} et donne
\[
 c_{\pi,n}(s)=\langle\pi(a_s)e_n,e_0\rangle,
 \qquad |c_{\pi,n}(s)|\le\Xi(a_s).
\]
Avec $\alpha=5/8>1/2$, Cauchy--Schwarz donne
\[
 \left|\sum_n f_{\pi,n}c_{\pi,n}(s)\right|
 \le C_\alpha\Xi(a_s)
       \left(\sum_n(1+n^2)^\alpha|f_{\pi,n}|^2\right)^{1/2},
 \quad C_\alpha^2=\sum_n(1+n^2)^{-\alpha}<\infty.
 \tag{S2}
\]
Cette étape n'utilise aucune borne sans norme de régularité pour un
vecteur arbitraire possédant une infinité de $K$-types.

Soit $u_{\pi,0}$ la réalisation du vecteur sphérique comme forme de
Maass ou série d'Eisenstein sur la surface modulaire, avec sa mesure
spectrale $d\nu(\pi)$. Pour $q=9/8>1$, le noyau ponctuel
\[
 J_i(q)=\int |u_{\pi,0}(i)|^2(1+\lambda_\pi)^{-q}\,d\nu(\pi)
 \tag{S3}
\]
est fini. L'intégrale comprend la somme cuspidale et l'intégrale
d'Eisenstein. On peut le voir directement par le calcul fonctionnel du
Laplacien de surface:
\[
 J_i(q)\le\frac1{\Gamma(q)}
     \int_0^\infty t^{q-1}e^{-t}H_t(i,i)\,dt.
\]
Le noyau de chaleur local satisfait $H_t(i,i)=O_i(t^{-1})$ pour
$0<t\le1$. Cette estimation reste vraie au point elliptique $i$:
son stabilisateur fini ne change que la constante. En effet le noyau
du quotient est la somme des noyaux hyperboliques sur l'orbite de $i$;
la croissance orbitale exponentielle et la décroissance gaussienne
du noyau pour $t\le1$ contrôlent cette somme. Pour $t\ge1$,
$H_t(i,i)\le H_1(i,i)$ par positivité de la mesure spectrale.
L'intégrale converge donc pour $q>1$.

Appliquons Cauchy--Schwarz une seconde fois, dans la variable spectrale.
En utilisant (S2)--(S3), nous obtenons
\[
 |\mathcal M_s f-\mu(f)|
 \le C_\alpha J_i(q)^{1/2}\Xi(a_s)
 \left(\int(1+\lambda_\pi)^q
                  \sum_n(1+n^2)^\alpha|f_{\pi,n}|^2\,d\nu(\pi)
 \right)^{1/2}.
\]
Sur le spectre sphérique non trivial, $\lambda_\pi\ge1/4$ et
\[
 (1+\lambda)^q(1+n^2)^\alpha
 \le C(1+\lambda+n^2/2)^{q+\alpha}.
\]
L'opérateur $\mathcal D$ agit sur ce type par $\lambda+n^2/2$.
Comme $q+\alpha=7/4$ et
$\Xi(a_s)=O((1+s)e^{-s})$, cela prouve (S1) pour des fonctions lisses
à support compact. Le passage à $W^{7/4}$ suit par densité et par le
théorème de trace local sur la courbe, de codimension~2; l'ordre
$7/4>1$ suffit pour cette trace. La preuve utilise seulement une
estimation locale au point de départ fixé $i$, et non une compacité
du quotient ni une borne uniforme du rayon d'injectivité dans la pointe.
\end{proof}

\section{Lissage uniforme sur le groupe}

Choisissons un noyau positif $\psi_h\in C_c^\infty(G)$ de masse~1,
supporté à distance $O(h)$ de l'identité, tel que
\[
 \|\psi_h\|_1=1,\qquad
 \|\mathcal L_Z\psi_h\|_1\le C_Zh^{-1}
\]
pour chacun des trois champs de multiplication à gauche sur $G$.
Posons $f_h=\int_G\psi_h(u)R(u)f\,du$.

\begin{lemma}[Normes du lissage]
Il existe $C$ indépendant de $0<h\le1$ tel que
\[
 \|f_h\|_{W^1}\le C\|f\|_{W^1},\qquad
 \|f_h\|_{W^2}\le Ch^{-1}\|f\|_{W^1},\qquad
 \|f_h\|_{W^{7/4}}\le Ch^{-3/4}\|f\|_{W^1}.
 \tag{S4}
\]
Si $f$ est globalement Lipschitz pour la métrique descendue d'une
métrique invariante à gauche sur $G$, alors
\[
 \|f_h-f\|_\infty\le Ch\operatorname{Lip}(f).
 \tag{S5}
\]
\end{lemma}

\begin{proof}
Notons $D_Z=dR(Z)$. L'identité
\[
 D_ZR(u)=R(u)D_{\operatorname{Ad}(u^{-1})Z}
\]
et le fait que $\operatorname{Ad}(u^{-1})$ reste borné sur le support des
noyaux donnent la première borne de (S4). Plus explicitement, écrivons
$\operatorname{Ad}(u^{-1})Z=\sum_j c_{Zj}(u)Z_j$. Alors
\[
 D_Z f_h=\sum_j\int_G\psi_h(u)c_{Zj}(u)R(u)D_{Z_j}f\,du.
\]
Une deuxième dérivée se déplace sur le noyau par intégration par parties
dans $u$:
\[
 D_YD_Z f_h
 =-\sum_j\int_G
    \mathcal L_Y\bigl(\psi_hc_{Zj}\bigr)(u)R(u)D_{Z_j}f\,du.
\]
Les fonctions $c_{Zj}$ et leurs premières dérivées sont uniformément
bornées sur un voisinage fixe de l'identité. Le membre de droite a
norme $L^2$ au plus $Ch^{-1}\|f\|_{W^1}$. Puisque
$\mathcal D=-\sum_jD_{Z_j}^2$, cela contrôle
$\|(1+\mathcal D)f_h\|_2=\|f_h\|_{W^2}$.
L'interpolation spectrale de l'opérateur positif $1+\mathcal D$ donne
\[
 \|f_h\|_{W^{7/4}}
 \le\|f_h\|_{W^1}^{1/4}\|f_h\|_{W^2}^{3/4},
\]
d'où la troisième borne de (S4). Ces calculs portent sur une
représentation unitaire du groupe; les constantes ne dépendent pas
des coordonnées de la pointe.

Pour (S5), la distance sur le quotient entre $\Gamma g$ et $\Gamma gu$
est majorée par la distance sur le groupe entre $g$ et $gu$, laquelle
vaut la distance entre $e$ et $u$ pour la métrique invariante à gauche.
Elle est donc $O(h)$ sur le support de $\psi_h$. Il suffit d'intégrer
l'inégalité de Lipschitz.
\end{proof}

\begin{theorem}[Critère de Lipschitz impliquant H]
Soit $f\in W^1(X)$ globalement Lipschitz et bornée. Alors
\[
 |\mathcal M_s f-\mu(f)|
 \le C(1+s)e^{-4s/7}
       \bigl(\|f\|_{W^1}+\operatorname{Lip}(f)\bigr).
 \tag{S6}
\]
En particulier, l'hypothèse H pour $f$ vaut avec tout
$0<\delta<1/14$.
\end{theorem}

\begin{proof}
Les deux moyennes sont des probabilités, donc l'erreur produite par le
lissage est au plus $2\|f_h-f\|_\infty$. Les propositions précédentes
fournissent
\[
 |\mathcal M_s f-\mu(f)|
 \le Ch\operatorname{Lip}(f)
      +C(1+s)e^{-s}h^{-3/4}\|f\|_{W^1}.
\]
Choisissons $h=e^{-4s/7}$. Les deux termes ont le taux exponentiel
$4/7$, ce qui prouve (S6). Puis
$4/7-(1/2+\delta)=1/14-\delta>0$ permet d'absorber le facteur $1+s$
dans la constante du $O(e^{-(1/2+\delta)s})$.
\end{proof}

\paragraph{Portée exacte.}
Ce théorème est une implication analytique. Son application à
$f=\Psi-E_g$ exige une preuve séparée de sa Lipschitzianité globale
et de sa norme $W^1$. Une fonction bornée globalement Lipschitz sur
ce quotient de volume fini appartient automatiquement à $W^1$:
son gradient faible est borné et la mesure riemannienne est un
multiple de Haar. Une preuve géométrique de cette Lipschitzianité
globale suffit donc, mais elle ne doit pas être remplacée par
l'estimation de valeur $|f|=O(\lambda_1)$ seule.

\paragraph{Sources et attribution.}
La borne par $K$-type est celle de Cowling--Haagerup--Howe, reproduite
dans Hongyu He, \emph{Bounds on smooth matrix coefficients on
$L^2$-spaces}, th.~3.1. L'absence d'exceptionnel de niveau~1 est rappelée
avec attribution à Maass/Roelcke dans E.~Yoshida, \emph{On the order
of growth of the Kloosterman zeta function}, J.~Math.~Soc.~Japan~44
(1992), remarque~1, p.~54. La normalisation
$\mathcal D=\mathcal C-2\Theta^2$ et les normes fractionnaires se trouvent
notamment dans E.~Corso et D.~Ravotti, \emph{Large hyperbolic circles},
arXiv:2208.07771, \S2.3. Leurs théorèmes d'équirépartition sont formulés
pour un quotient compact et ne sont pas utilisés ici comme théorèmes
sur la surface modulaire. La combinaison par les deux applications de
Cauchy--Schwarz, le noyau ponctuel et le lissage de groupe est explicitée
ci-dessus pour le quotient non compact, spectre continu compris.
